% =====================================================================
%  accessible-LaTeX-template.tex
%  Template: PDF/UA-2 tagged PDF with MathML for course handouts
%
%  HOW TO USE
%    1. Fill in the handout details in the "EDIT THESE" block bealow.
%    2. Replace the placeholder sections, text, equations, table,
%       and figure with your own content.
%    3. Compile with LuaLaTeX (required for MathML generation).
%    4. Check the result with veraPDF (PDF/UA-2 profile) and, if you
%       can, listen to it with a screen reader.
%
%  More information available at https://kb.swarthmore.edu/wiki/Accessible_LaTeX
%  Requires a recent TeX distribution (TeX Live 2025 or newer).
%
%  Lines marked [A11Y] indicate accessibility-related settings.
% =====================================================================

% [A11Y] Turns on tagging, sets the document language, and makes the
% output a PDF/UA-2 file. Change "lang" if the handout is not in
% US English (e.g. en-GB, es, fr).
\DocumentMetadata{
  lang          = en-US,
  pdfversion    = 2.0,
  pdfstandard   = ua-2,
  tagging       = on,
  tagging-setup = {math/setup = mathml-SE} % [A11Y] math as MathML
}

% [A11Y] 12pt text is easier to read on paper, where readers can't zoom.
\documentclass[12pt]{article}

% ---------------------------------------------------------------------
%  EDIT THESE: handout details (used for both the visible title and
%  the PDF metadata, so they always match)
% ---------------------------------------------------------------------
\newcommand{\HandoutTitle}{Handout Title}
\newcommand{\HandoutAuthor}{Instructor Name}
\newcommand{\HandoutCourse}{Course Number: Course Name}
\newcommand{\HandoutDate}{Term Year}

% ---------------------------------------------------------------------
%  Packages
%  Order matters: amsmath before unicode-math, hyperref near the end.
% ---------------------------------------------------------------------
\usepackage[margin=1in]{geometry}
\usepackage{amsmath}
\usepackage{amsthm}
\usepackage{unicode-math}           % OpenType math font with Unicode mapping
\setmainfont{Latin Modern Roman}    % change the text font here
\setmathfont{Latin Modern Math}     % change the math font here
\usepackage{graphicx}
\usepackage{booktabs}               % clean table rules
\usepackage{hyperref}

% [A11Y] Title metadata is required by PDF/UA. Links are shown as
% colored text with strong contrast instead of boxes.
\hypersetup{
  pdftitle           = {\HandoutTitle},
  pdfauthor          = {\HandoutAuthor},
  pdfdisplaydoctitle = true,
  colorlinks         = true,
  linkcolor          = blue!60!black,
  urlcolor           = blue!60!black,
  citecolor          = blue!60!black
}

% ---------------------------------------------------------------------
%  Theorem-style environments (rename or remove as needed)
% ---------------------------------------------------------------------
\theoremstyle{definition}
\newtheorem{definition}{Definition}[section]
\newtheorem{example}[definition]{Example}
\theoremstyle{plain}
\newtheorem{theorem}[definition]{Theorem}

% ---------------------------------------------------------------------
\title{\HandoutTitle}
\author{\HandoutAuthor\\ \HandoutCourse}
\date{\HandoutDate}
% ---------------------------------------------------------------------

\begin{document}

% [A11Y] The title becomes the top-level heading; sections sit under it.
\maketitle

% =====================================================================
\section{Section Title}
% =====================================================================

Introductory text for this section. Briefly state what the reader will
learn and define any notation before using it, so that someone
listening with a screen reader has the context they need.

\begin{definition}
  State a definition here. Use \emph{emphasis} for the term being
  defined.
\end{definition}

\subsection{Subsection Title}

Text for this subsection. A displayed, numbered equation looks like
this (replace the placeholder with your own):
\begin{equation}\label{eq:first}
  y = f(x)
\end{equation}

% [A11Y] Refer to equations by label so readers get a working link.
Refer back to it as Equation~\eqref{eq:first}. Inline math is written
like this: \(f(x)\).

\subsection{Subsection Title}

For multi-line work, use \texttt{align}:
\begin{align}
  \text{left side} &= \text{first step} \label{eq:step1} \\
                   &= \text{second step} \label{eq:step2}
\end{align}

A bulleted list:
\begin{itemize}
  \item First point.
  \item Second point.
  \item Third point.
\end{itemize}

% =====================================================================
\section{Section Title}
% =====================================================================

\begin{theorem}
  State a theorem here.
\end{theorem}

\begin{proof}
  Write the proof here.
\end{proof}

\begin{example}
  Work through an example here.
\end{example}

\subsection{Table}

Table~\ref{tab:placeholder} shows how to lay out a table.
% [A11Y] Introduce every table in the text before it appears.

\begin{table}[htbp]
  % [A11Y] Marks the first row as column headers for screen readers.
  % Change the number if your table has more than one header row.
  \tagpdfsetup{table/header-rows={1}}
  \centering
  % [A11Y] Captions go above tables and should say what the table shows.
  \caption{Describe what this table shows}
  \label{tab:placeholder}
  % [A11Y] Keep tables simple: one header row, no merged cells, and
  % narrow enough to fit the page. Put math in \( \) in every cell,
  % headers included.
  \begin{tabular}{lccc}
    \toprule
    Header 1 & Header 2 & Header 3 & Header 4 \\
    \midrule
    Row label & Entry & Entry & Entry \\
    Row label & Entry & Entry & Entry \\
    Row label & Entry & Entry & Entry \\
    \bottomrule
  \end{tabular}
\end{table}

\subsection{Figure}

Figure~\ref{fig:placeholder} shows how to include an image.

\begin{figure}[htbp]
  \centering
  % [A11Y] Every image needs alt text. Describe what the reader should
  % take away, not just what the image is. For a graph, say
  % something like: "Parabola opening upward, crossing the x-axis at
  % x = 1 and x = 2, with its lowest point at x = 1.5."
  % "example-image" is a placeholder bundled with TeX Live; replace
  % it with your own file name.
  \includegraphics[width=0.6\textwidth,
                   alt={Replace with a description of the image}]
                  {example-image}
  \caption{Describe what this figure shows}
  \label{fig:placeholder}
\end{figure}

% =====================================================================
\section{Exercises}
% =====================================================================

\begin{enumerate}
  \item First exercise.
  \item Second exercise.
  \item Third exercise.
\end{enumerate}

\end{document}